\documentclass[12pt,letterpaper]{article}
\usepackage[margin=0.8in]{geometry}
\usepackage{amsmath,amssymb}
\setlength{\parindent}{0pt}
\setlength{\parskip}{0.7em}
\pagestyle{plain}
\newcommand{\question}[1]{\noindent{\large\bfseries #1.}\hfill{\small Generated}\par\medskip}
\newcommand{\partlabel}[1]{\noindent\textbf{(#1)}\ }
\begin{document}
\begin{titlepage}
\thispagestyle{empty}
\vspace*{0.6in}
{\Large PHYS 200\par}
\vspace{0.25in}
{\Huge\bfseries Practice Set 3\par}
\vspace{0.18in}
{\large 12 questions\par}
\vspace{0.4in}
\hrule
\vspace{0.25in}
{\Large\bfseries Time limit: 240 minutes (4 hours)\par}
This is a generous practice limit, already including your double-time allowance.
Allow about 180--210 minutes for the questions and use the remaining time to
finish or check your work. This is a practice estimate, not an official exam limit.
\vspace{0.2in}

\textbf{Instructions}
\begin{itemize}
\setlength{\itemsep}{0.55em}
\item Answer all 12 questions. Each question starts on a new page with space for working.
\item Numerical arithmetic is not required. A correct symbolic expression, or
a sufficient set of equations determining the requested quantity, is a complete
answer. Define any additional unknowns and state any necessary physical conditions.
\item Show enough reasoning to justify your setup and conclusions. A named
method alone is not sufficient.
\item Identify the frames and direction conventions in your work. Define any
Lorentz factor you use, retain signs, and keep units consistent.
\item Use your usual course formula sheet if desired. No lengthy derivations
are required.
\end{itemize}

\vfill
Name:\ \rule{3.6in}{0.4pt}

Date:\ \rule{1.6in}{0.4pt}
\hfill Start:\ \rule{1.1in}{0.4pt}
\quad Finish:\ \rule{1.1in}{0.4pt}
\vspace*{0.2in}
\end{titlepage}
\setcounter{page}{1}

\question{1}
A stationary bound system consists of two constituents, each of rest mass $m$.
Its positive binding energy $B$ is unknown. It decays into exactly two particles,
each of rest mass $\mu$, where $0<\mu<m$. In the original system's rest frame,
each outgoing particle is measured to have speed $w$, with $0<w<c$.
There are no other products.

\partlabel{a} Determine $B$ in terms of the measured speed and the given masses.
\vfill

\partlabel{b} What condition must the measured speed satisfy for the inferred
binding energy to be positive?
\vfill
\newpage

\question{2}
A metal is illuminated first with monochromatic light of unknown frequency $f$,
then with light of frequency $2f$. The measured stopping-potential magnitudes
are $V_1$ and $V_2$, respectively, with $V_2>2V_1>0$.
The work function is the same in both measurements. Let $e$ denote the positive
magnitude of the electron charge.

\partlabel{a} Determine the frequency $f$ and the metal's work function from
these measurements.
\vfill

\partlabel{b} The first light source is used again at frequency $f$, but its
intensity is trebled over the same illuminated area. What happens to the maximum
kinetic energy of the emitted electrons, the stopping-potential magnitude, and
the electron emission rate? Assume constant emission efficiency and no saturation.
\vfill
\newpage

\question{3}
Two photons with unequal energies collide in the laboratory and produce a
single particle of measured rest mass $M$. There are no other products.
One incoming photon has known energy $\mathcal E$ and travels along $+x$.
The other travels in the $xy$-plane at $120^\circ$ counterclockwise from $+x$;
its energy is unknown.

\partlabel{a} Determine the unknown incoming photon energy.
\vfill

\partlabel{b} Determine both laboratory velocity components of the product,
including their signs.
\vfill
\newpage

\question{4}
In an inertial frame $S$, event A occurs at $x=0$, $t=0$, and event B occurs at
$x=-L$, $t=T$, where $L>cT>0$. An observer moving along the $x$-axis says that
the two events are simultaneous in their own inertial frame.

\partlabel{a} Find the observer's signed velocity relative to $S$ and the signed
spatial separation $x'_B-x'_A$ in that observer's frame.
\vfill

\partlabel{b} Could a different inertial observer find that the two events
occurred at the same location? Justify your answer.
\vfill
\newpage

\question{5}
A particle of rest mass $M$ is initially at rest. It decays into one particle of
rest mass $m$ and two photons, with $M>m>0$. No other particles are produced.
Treat this as a question of kinematic possibility only.

What is the largest possible speed of the massive daughter in the original
parent's rest frame? Describe the directions of the two photons relative to
the daughter when that largest speed is attained, and justify why a larger
speed is impossible.

A symbolic expression or equations determining the largest speed are sufficient.
\vfill
\newpage

\question{6}
A receiver remains at the origin of an inertial frame $S$. A source moves with
constant velocity $-v\,\hat{\mathbf x}$ in that frame, where $0<v<c$.
At the instant it emits a photon that will reach the receiver, the source is at
$(a,b)$ in $S$, with $a>0$ and $b>0$. The emitted light has wavelength
$\lambda_0$ in the source's rest frame. Light travels through vacuum.

\partlabel{a} Find the wavelength measured by the receiver.
\vfill

\partlabel{b} At emission the source's distance from the receiver is decreasing.
Can the detected wavelength nevertheless be longer than $\lambda_0$?
Give the condition on $a$, $b$, and $v$ that decides this.
\vfill
\newpage

\question{7}
A particle of rest mass $5m$ moves at velocity $-V\,\hat{\mathbf x}$ in the
laboratory, where $0<V<c$. It decays into daughter A of rest mass $m$ and
daughter B of rest mass $2m$. In the parent's rest frame, A moves along $+x$
and B along $-x$. No other particles are produced.

\partlabel{a} Find the total energy and signed $x$-momentum of each daughter
in the parent's rest frame.
\vfill

\partlabel{b} Find the total energy and signed $x$-momentum of each daughter
in the laboratory.
\vfill

\partlabel{c} Find A's signed laboratory velocity. For what range of parent
speeds $V$ does A move along $-x$ in the laboratory?
\vfill
\newpage

\question{8}
A photon scatters from a free electron initially at rest in the laboratory.
A detector measures the scattered photon's wavelength $\Lambda$ and the
electron's kinetic energy $K>0$. The incident wavelength and the photon
scattering angle were not measured. The electron has rest mass $m_e$.

\partlabel{a} Determine the incident wavelength and the angle between the
incident and scattered photon directions.
\vfill

\partlabel{b} Can every positive pair of measured values $(\Lambda,K)$ arise
in this experiment? State a condition on the measurements that decides whether
they are physically consistent.
\vfill
\newpage

\question{9}
Two particles, each of rest mass $m$, collide and stick together to make one
object, with no radiation or other products. Their individual total laboratory
energies are known to be $E_1$ and $E_2$, with $E_1>E_2>mc^2$.
Their incoming directions were not recorded.

\partlabel{a} Do these measurements determine a unique rest mass for the
product? Explain.
\vfill

\partlabel{b} Find the smallest and largest product rest masses consistent
with the measurements. Identify a physically possible incoming geometry for
each extreme.
\vfill
\newpage

\question{10}
A spacecraft moves horizontally with velocity $-V\,\hat{\mathbf x}$ relative
to the ground, where $0<V<c$. It launches a small probe. Ground observers
measure the probe to be moving straight upward at speed $w$, where $0<w<c$.

\partlabel{a} Find the probe's two velocity components in the spacecraft's
rest frame and its launch angle measured counterclockwise from that frame's
$+x$ direction. The two frames' axes are parallel.
\vfill

\partlabel{b} Use your result to show that the spacecraft also measures a
probe speed less than $c$.
\vfill
\newpage

\question{11}
In frame $S$, one photon travels along $+x$ and another along $-x$.
Their combined energy in $S$ is $U$, but their individual energies are unknown.
The inertial frame $Z$ in which their combined momentum vanishes moves at
velocity $-u\,\hat{\mathbf x}$ relative to $S$, where $0<u<c$.

\partlabel{a} Determine the energy of each photon in $S$, identifying which
belongs to each direction.
\vfill

\partlabel{b} Determine the invariant mass of the two-photon system and the
energy of each photon measured in $Z$.
\vfill
\newpage

\question{12}
An isolated body of initial rest mass $m$ is at rest in the laboratory.
It absorbs a photon of laboratory energy $\varepsilon>0$ travelling along
$-x$. Afterward there is just one body, with the absorbed energy retained as
internal energy and motion; nothing is emitted.

\partlabel{a} Determine the final body's rest mass and signed laboratory
velocity.
\vfill

\partlabel{b} A proposed model of this same process assumes that the final
body still has rest mass $m$. Is that model consistent with the stated
experiment? Justify your answer.
\vfill

\end{document}
